% Part: many-valued-logic
% Chapter: three-valued-logics
% Section: lukasiewicz

\documentclass[../../../include/open-logic-section]{subfiles}

\begin{document}

\olfileid{mvl}{thr}{luk}

\olsection{منطق \L ukasiewicz}

یکی از نخستین طرح‌های منتشرشده و پرورده برای یک منطق چندارزشی را
فیلسوف لهستانی، یان \L ukasiewicz، در سال~۱۹۲۱ ارائه کرد. انگیزهٔ
\L ukasiewicz مسئلهٔ نبرد دریایی ارسطو بود: به نظر می‌رسد جملهٔ «فردا
نبردی دریایی رخ خواهد داد» \emph{امروز} نه صادق است و نه کاذب، زیرا
ارزش صدقش هنوز تعیین نشده است. \L ukasiewicz پیشنهاد کرد ارزش صدق
سومی را برای چنین جمله‌های «ممکنِ آینده» در نظر بگیریم.
\begin{quote}
می‌توانم بی‌هیچ تناقضی فرض کنم که حضور من در ورشو در لحظه‌ای معین از
سال آینده، مثلاً ظهر ۲۱ دسامبر، در زمان حاضر نه به‌صورت مثبت و نه
به‌صورت منفی تعیین شده است. بنابراین ممکن، اما نه ضروری، است که در آن
زمان معین در ورشو حضور داشته باشم. بنا بر این فرض، گزارهٔ «ظهر ۲۱
دسامبر سال آینده در ورشو خواهم بود» در زمان حاضر نمی‌تواند نه صادق
باشد و نه کاذب. زیرا اگر اکنون صادق بود، حضور آیندهٔ من در ورشو باید
ضروری می‌بود، که با فرض تناقض دارد. از سوی دیگر، اگر اکنون کاذب بود،
حضور آیندهٔ من در ورشو باید ناممکن می‌بود، که آن نیز با فرض تناقض
دارد. پس گزارهٔ موردنظر در این لحظه نه صادق است و نه کاذب و باید ارزش
سومی، متفاوت از «۰» یا کذب و «۱» یا صدق، داشته باشد. می‌توانیم این
ارزش را با «$\frac{1}{2}$» نشان دهیم. این ارزش نمایندهٔ «ممکن» است و
به‌عنوان ارزش سوم به «صادق» و «کاذب» می‌پیوندد.
\end{quote}
برای ارزش صدق سوم \L ukasiewicz از~$\Undef$ استفاده می‌کنیم.
\footnote{\L ukasiewicz در اینجا «ممکن» را به معنایی به کار می‌برد که
امروزه رایج نیست: ممکن اما نه ضروری.}

توابع صدق رابط‌های~$\lnot$، $\land$ و~$\lor$ در این تفسیر به‌آسانی
تعیین می‌شوند: نقیض یک جملهٔ ممکنِ آینده نیز جمله‌ای ممکنِ آینده است؛
پس~$\tf{\lnot}(\Undef) = \Undef$. اگر یک مؤلفهٔ عطف نامتعین و دیگری
صادق باشد، عطف نیز نامتعین است؛ زیرا بسته به اینکه مؤلفهٔ ممکنِ آینده
چه سرانجامی پیدا کند، عطف ممکن است صادق یا کاذب شود. پس
\[
    \tf{\land}(\True, \Undef) =
\tf{\land}(\Undef, \True) =
\Undef.
\]
اما اگر مؤلفهٔ دیگر کاذب باشد، عطف نمی‌تواند صادق شود؛ پس
\[\tf{\land}(\False, \Undef) =
\tf{\land}(\Undef, \False) = \False.\]
دیگر مقادیر (اگر آرگومان‌ها ارزش‌های صدق تعیین‌شده، یعنی~$\True$ یا
$\False$، باشند) مانند منطق کلاسیک‌اند.

وضع رابط شرطی کمی پیچیده‌تر است. فرض کنید~$q$ گزاره‌ای ممکنِ آینده
باشد. اگر~$p$ کاذب باشد، $p \lif q$ مستقل از سرانجام~$q$ صادق خواهد
بود؛ پس باید~$\tf{\lif}(\False, \Undef) = \True$ بگذاریم. و اگر~$p$
صادق باشد، $q \lif p$ مستقل از سرانجام~$q$ صادق خواهد بود؛ پس
$\tf{\lif}(\Undef, \True) = \True$. اگر~$p$ صادق باشد، $p \lif q$
ممکن است سرانجام صادق یا کاذب شود؛ پس
$\tf{\lif}(\True, \Undef) = \Undef$. به‌همین‌ترتیب، اگر~$p$ کاذب
باشد، $q \lif p$ ممکن است سرانجام صادق یا کاذب شود؛ پس
$\tf{\lif}(\Undef, \False) = \Undef$. تنها حالتی باقی می‌ماند که~$p$
و~$q$ هر دو ممکنِ آینده باشند. بنا بر انگیزهٔ اصلی، در این حالت واقعاً
باید~$\Undef$ را نسبت دهیم. اما در آن صورت~$!A \lif !A$ یک همان‌گویی
\emph{نخواهد بود}. \L ukasiewicz از کنارگذاشتن~$!A \lor \lnot !A$ و
$\lnot(!A \land \lnot !A)$ ابایی نداشت، اما از کنارگذاشتن
$!A \lif !A$ پرهیز کرد. ازاین‌رو مقرر کرد
$\tf{\lif}(\Undef, \Undef) = \True$.

\begin{defn}\ollabel{def:lukasiewicz}
منطق سه‌ارزشی \L ukasiewicz با ماتریس زیر تعریف می‌شود:
\begin{enumerate}
  \item زبان گزاره‌ای استاندارد~$\Lang L_0$ با~$\lnot$، $\land$،
  $\lor$ و~$\lif$.
  \item مجموعهٔ ارزش‌های صدق~$V = \{\True, \Undef, \False\}$.
  \item $\True$ تنها ارزش ممتاز است؛ یعنی~$V^+ = \{\True\}$.
  \item توابع صدق با جدول‌های زیر داده می‌شوند:
  \begin{center}
    \begin{tabular}{c|c} 
      $\tf{\lnot}$ & \\ 
      \hline  
      $\True$ & $\False$ \\ 
      $\Undef$ & $\Undef$ \\
      $\False$ & $\True$ 
    \end{tabular}
    \quad
    \begin{tabular}{c|ccc} 
      $\tf{\land}[\LogLuk[3]]$ & $\True$ & $\Undef$ & $\False$ \\ 
      \hline 
      $\True$ & $\True$ & $\Undef$ & $\False$ \\ 
      $\Undef$ & $\Undef$ & $\Undef$ & $\False$\\ 
      $\False$ & $\False$ & $\False$ & $\False$ 
    \end{tabular}
    \\[2ex]
    \begin{tabular}{c|ccc} 
      $\tf{\lor}[\LogLuk[3]]$ & $\True$ & $\Undef$ & $\False$ \\ 
      \hline 
      $\True$ & $\True$ & $\True$ & $\True$ \\ 
      $\Undef$ & $\True$ & $\Undef$ & $\Undef$ \\
      $\False$ & $\True$ & $\Undef$ & $\False$ 
    \end{tabular}
    \quad
    \begin{tabular}{c|ccc} 
      $\tf{\lif}[\LogLuk[3]]$ & $\True$ & $\Undef$ & $\False$ \\ 
      \hline 
      $\True$ & $\True$ & $\Undef$ & $\False$ \\ 
      $\Undef$ & $\True$ & $\True$ & $\Undef$  \\ 
      $\False$ & $\True$ & $\True$ & $\True$ 
    \end{tabular}
  \end{center} 
\end{enumerate}
\end{defn}

به‌آسانی می‌توان دید هر~!!{formula} چون~$!A$ که فقط شامل~$\lnot$،
$\land$ و~$\lor$ باشد، اگر به همهٔ~!!{propositional variable}هایش
$\Undef$ نسبت داده شود، ارزش صدق~$\Undef$ می‌گیرد. پس برای نمونه،
همان‌گویی‌های کلاسیک~$p \lor \lnot p$ و~$\lnot(p \land \lnot p)$ در
$\LogLuk[3]$ همان‌گویی نیستند، زیرا هرگاه~$\pAssign v(p) = \Undef$،
$\pValue{v}(!A) = \Undef$.

در~!!{valuation}هایی که~$\pAssign v(p)$ برابر~$\True$ یا~$\False$ است،
$\pValue v(!A)$ با ارزش صدق کلاسیک آن یکسان خواهد بود.

\begin{prop}
  اگر برای همهٔ~$p$های موجود در~$!A$،
  $\pAssign v(p) \in \{\True, \False\}$، آنگاه
  $\pValue v[\LogLuk[3]](!A) = \pValue v[\LogCL](!A)$.
\end{prop}

\begin{prob}\label{mvl:thr:luk:prob:luk-iff} فرض کنید در~$\LogLuk[3]$
  تعریف کنیم
  $\pValue v(!A \liff !B) = \pValue v((!A \lif !B) \land (!B \lif
  !A))$. در آن صورت~$\liff$ چه جدول صدقی خواهد داشت؟
\end{prob}

بسیاری از همان‌گویی‌های کلاسیک در~\LogLuk[3] نیز همان‌گویی‌اند؛ برای
نمونه~$\lnot p \lif (p \lif q)$. درست مانند منطق کلاسیک، می‌توانیم این
امر را با جدول صدق وارسی کنیم:
\begin{center}
\begin{tabular}{cc|cccccc}
  $p$ & $q$ & $\lnot$ & $p$ & $\lif$ & $\smash{(}p$ & $\lif$ & $q\smash{)}$ \\ \hline
  \True & \True & \False & \True & \True & \True & \True & \True \\
  \True & \Undef & \False & \True & \True & \True & \Undef & \Undef \\
  \True & \False & \False & \True & \True & \True & \False & \False \\
  \Undef & \True & \Undef & \Undef & \True & \Undef & \True & \True \\
  \Undef & \Undef & \Undef & \Undef & \True & \Undef & \True & \Undef \\
  \Undef & \False & \Undef & \Undef & \True & \Undef & \Undef & \False \\
  \False & \True & \True & \False & \True & \False & \True & \True \\
  \False & \Undef & \True & \False & \True & \False & \True & \Undef \\
  \False & \False & \True & \False & \True & \False & \True & \False \\  
\end{tabular}
\end{center}

\begin{prob}
  نشان دهید عبارت‌های زیر در~\LogLuk[3] همان‌گویی‌اند:
  \begin{enumerate}
    \item $p \lif (q \lif p)$
    \item\label{mvl:thr:luk:prob:luk-taut-2} $\lnot(p \land q) \liff (\lnot p \lor \lnot q)$
    \item\label{mvl:thr:luk:prob:luk-taut-3} $\lnot(p \lor q) \liff (\lnot p \land \lnot q)$
  \end{enumerate}
  (در~\olref[mvl][thr][luk]{prob:luk-taut-2} و
  \olref[mvl][thr][luk]{prob:luk-taut-3}، عبارت~$!A \liff !B$ را مخفف
  $(!A \lif !B) \land (!B \lif !A)$ بگیرید، یا به پاسخ خود برای
  \olref[mvl][thr][luk]{prob:luk-iff} مراجعه کنید.)
\end{prob}

\begin{prob}
  نشان دهید همان‌گویی‌های کلاسیک زیر در~\LogLuk[3] همان‌گویی نیستند:
  \begin{enumerate}
    \item $(\lnot p \land p) \lif q$
    \item $((p \lif q) \lif p) \lif p$
    \item $(p \lif (p \lif q)) \lif (p \lif q)$
  \end{enumerate}
\end{prob}

شاید ازاین‌رو گمان شود که هرچند همهٔ همان‌گویی‌های کلاسیک در
$\LogLuk[3]$ همان‌گویی نیستند، باید دست‌کم در هر~!!{valuation} یکی از
دو مقدار~$\True$ یا~$\Undef$ را بگیرند. چنین نیست. مثال نقض عبارت زیر
است:
\[
  \lnot(p \lif \lnot p) \lor \lnot(\lnot p \lif p)
\]
که اگر~$p$ برابر~$\Undef$ باشد، مقدارش~$\False$ است.

\begin{prob}
  کدام‌یک از روابط زیر در منطق \L ukasiewicz برقرارند؟ برای هریک جدول
  صدق بدهید.
  \begin{enumerate}
    \item $p, p \lif q \Entails q$
    \item $\lnot\lnot p \Entails p$
    \item $p \land q \Entails p$
    \item $p \Entails p \land p$
    \item $p \Entails p \lor q$
  \end{enumerate}
\end{prob}

\L ukasiewicz امیدوار بود با افزودن رابط یک‌موضعی~$\Diamond !A$ (برای
«$!A$ ممکن است») و رابط متناظر~$\Box !A$ (برای «$!A$ ضروری است») بر
پایهٔ دستگاه سه‌ارزشی خود منطقی برای امکان بسازد:
\begin{center}
  \begin{tabular}{c|c} 
    $\tf{\Diamond}$ & \\ 
    \hline  
    $\True$ & $\True$ \\ 
    $\Undef$ & $\True$ \\
    $\False$ & $\False$ 
  \end{tabular}
  \quad
  \begin{tabular}{c|c} 
    $\tf{\Box}$ & \\ 
    \hline  
    $\True$ & $\True$ \\ 
    $\Undef$ & $\False$ \\
    $\False$ & $\False$ 
  \end{tabular}
\end{center}
به بیان دیگر، $p$ ممکن است اگر و تنها اگر از پیش به‌عنوان کاذب تعیین
نشده باشد؛ و~$p$ ضروری است اگر و تنها اگر از پیش به‌عنوان صادق تعیین
شده باشد.

\begin{prob}
  نشان دهید~$\Box p \liff \lnot\Diamond \lnot p$ و
  $\Diamond p \liff \lnot \Box \lnot p$ در~\LogLuk[3]، پس از افزودن
  جدول‌های صدق~$\Box$ و~$\Diamond$، همان‌گویی‌اند.
\end{prob}

بااین‌حال، کاستی‌های این منطق موجهات پیشنهادی زود آشکار شد: رویدادها
هرگونه که پیش بروند، $p \land \lnot p$ هرگز نمی‌تواند صادق شود. پس حتی
اگر اکنون تعیین‌نشده و در نتیجه نامتعین باشد، باید ناممکن شمرده شود؛
یعنی~$\lnot \Diamond(p \land \lnot p)$ باید همان‌گویی باشد. اما اگر
$\pAssign v(p) = \Undef$، آنگاه
$\pValue v(\lnot \Diamond(p \land \lnot p)) = \Undef$. هرچند
\L ukasiewicz درست می‌گفت که دو ارزش صدق برای تمایزهای موجهاتی مانند
امکان و ضرورت کافی نیست، افزودن ارزش صدق سوم نیز کافی نیست.

\end{document}
